\section{Related Work}
\label{sec:related}

\subsection{Streaming and seasonal anomaly detection}
Detecting departures in a running time series is well studied. Control charts based
on exponentially weighted moving averages~\citep{roberts1959ewma} and rolling
$z$-scores are the classic lightweight approach, and our breach test is in this
family. Richer methods model seasonality and trend: Prophet decomposes a series
into trend, seasonality, and holidays~\citep{taylor2018prophet}, RobustSTL gives a
robust seasonal-trend decomposition for long series~\citep{wen2019robuststl}, and
seasonal hybrid ESD detects anomalies on top of such a decomposition~%
\citep{hochenbaum2017shesd}. Numenta's hierarchical temporal memory targets
unsupervised real-time detection on streams~\citep{ahmad2017htm}. Our detector
sits at the simple end of this spectrum on purpose. The contribution is not a new
detector but where it runs and how cheaply: a per-series $z$-score with
confirmation, expressed as a causaloid, light enough to run on every agent. We have
not compared detection quality against these methods, which \cref{sec:discussion}
names as the main open gap.

\subsection{Computational versus classical causality}
\dc{} implements the Effect Propagation Process, a computational model of
causality~\citep{deepcausality}. This is a different tradition from the
inference-focused causal models a reader may expect: Pearl's structural causal
models reason about interventions and counterfactuals within an assumed
structure~\citep{pearl2009causality}, and Granger causality tests predictive
precedence between series~\citep{granger1969causality}. We use the
computational-causality framing for its engineering properties, namely a small,
composable, testable unit of inference evaluated over an explicit context. We make
no causal-inference claim about the metrics themselves.
